\documentclass[UTF8,a4paper,11pt]{ctexart}
\usepackage[a4paper,margin=2cm,headheight=14pt]{geometry}
\usepackage{amsmath,amssymb,mathtools,bm}
\usepackage{booktabs,longtable,array,tabularx}
\usepackage{xcolor,xurl,hyperref,enumitem,microtype}
\hypersetup{colorlinks=true,linkcolor=blue,urlcolor=blue,citecolor=blue}
\urlstyle{tt}\Urlmuskip=0mu plus 3mu
\newcommand{\pathref}[1]{\begingroup\Urlmuskip=0mu plus 3mu\path{#1}\endgroup}
\newcommand{\code}[1]{\texttt{\detokenize{#1}}}
\newcolumntype{L}[1]{>{\raggedright\arraybackslash}p{#1}}
\newcolumntype{Y}{>{\raggedright\arraybackslash}X}
\newcommand{\C}{\mathbb C}\newcommand{\norm}[1]{\left\lVert#1\right\rVert}
\newcommand{\ip}[2]{\left\langle#1,#2\right\rangle}
\newcommand{\Yhat}{\widehat Y}\newcommand{\src}[1]{\par\smallskip\noindent\textit{参考源：}#1\par\smallskip}
\setlength{\parindent}{2em}\setlength{\parskip}{.35em}\emergencystretch=3em
\title{QUDA MultiGrid 粗层 $X/Y/\Yhat$、Clover-like dslash 与矩阵顺序\\
\large 从聚合基、Galerkin 投影到递归求解的完全构建文档}
\author{PyQCU 工程分析记录}\date{2026 年 9 月 13 日}
\begin{document}\maketitle
\begin{abstract}
本文是对 \pathref{refer/git-rep/quda} 中 Wilson/Clover MultiGrid 粗层实现的源码级补充解析，专门说明 onsite 矩阵 $X$、方向链接 $Y$、预条件链接 $\Yhat$ 的数学定义、存储位置、生成次序、Clover-like dslash 的真实行为以及矩阵左右顺序。QUDA 源码是第一证据，PyQCU strict full-coarse 实现是交叉验证。本文把 full operator、Clover-PC、coarse-PC 和 parity Schur 分开，最后以一个连续的超长单列表格给出从 null vector 到 $\Yhat$、MATPC 与验证的整条算法。
\end{abstract}
\tableofcontents

\section{范围、证据与六条结论}
本文只研究给定 fine Wilson/Clover operator、聚合 null vectors 和 coarse geometry 后，QUDA 如何生成并应用
\[
\{X_\ell,Y_{\ell,\mu}^{\rm f},Y_{\ell,\mu}^{\rm b},
X_\ell^{-1},\Yhat_{\ell,\mu}^{\rm f},\Yhat_{\ell,\mu}^{\rm b}\}.
\]
“构建”包括 field 分配、parity/方向映射、局部矩阵乘、atomic 累加、固定点缩放和 ghost exchange；“应用”包括 raw dslash、onsite block、preconditioned hopping、MATPC、prepare/reconstruct 及验证。

\begin{longtable}{@{}L{.18\textwidth}L{.30\textwidth}L{.42\textwidth}@{}}
\caption{证据分级}\label{tab:evidence}\\\toprule
等级&材料&用法\\\midrule\endfirsthead\toprule
等级&材料&用法\\\midrule\endhead
源码事实&QUDA QUDA coarse operator、dslash 和 preconditioned kernel&描述真实参数、分支、存储和生命周期。\\
数学等价&$R=P^\dagger$、$D_{\ell+1}=R A_\ell^{\rm eff}P$、Schur 消元、$\Yhat$ 定义&解释源码次序并做维数、共轭和量纲检查。\\
交叉实现&\pathref{pyqcu/tools/_strict_galerkin.py}、\pathref{pyqcu/solver/_quda_multigrid.py}、\pathref{cpp/cuda/qcu/src/apply_multigrid_strict.cu}&核对 full coarse、compact parity 和 QCU asset ABI。\\
未验证&尚未运行的 GPU 逐元素对照&明确标出，不能用推导代替数值证据。\\\bottomrule
\end{longtable}

结论先写清楚。$X$ 是 coarse site 上的 dense $E\times E$ 矩阵，不是 $SU(3)$；$Y_\mu^{\rm f}(X)$ 存在源点 $X$，乘 $z(X+\hat\mu)$；$Y_\mu^{\rm b}(X-\hat\mu)$ 存在 $X-\hat\mu$，action 读取其 dagger。raw action 为
\[
(D_cz)(X)=X(X)z(X)+\sum_\mu\{Y_\mu^{\rm f}(X)z(X+\hat\mu)+Y_\mu^{\rm b}(X-\hat\mu)^\dagger z(X-\hat\mu)\}.
\]
预条件链接必为
\[
\Yhat_\mu^{\rm f}(X)=X^{-1}(X)Y_\mu^{\rm f}(X),\qquad
\Yhat_\mu^{\rm b}(X-\hat\mu)=Y_\mu^{\rm b}(X-\hat\mu)X^{-\dagger}(X).
\]
\code{DiracCoarsePC::Dslash} 只应用 $\Yhat$ hopping；它传入 $X$ 但关闭 onsite 分支。Clover-PC 和 coarse-PC 需要双向链接；普通 Wilson-like setup 可按 spin projector 规则由一侧 reverse。
\src{\pathref{docs/report_multigrid_quda_pyqcu_20260909.tex:547--691} 是原报告对应章节；本文逐项展开其结论。}

\section{自由度、几何和量纲}
fine lattice 为 $\Lambda_f=(L_x,L_y,L_z,L_t)$，Wilson/Clover 每站点自由度为
\[
e=N_s^fN_c^f=4\cdot3=12.
\]
block $b=(b_x,b_y,b_z,b_t)$ 给出
\[
X_\mu=\lfloor x_\mu/b_\mu\rfloor,\quad
\Lambda_c=(L_x/b_x,L_y/b_y,L_z/b_z,L_t/b_t).
\]
coarse spin 为 $N_s^c=2$，null-vector 数为 $n_v$，故 $E=N_s^cn_v$，$X,Y,\Yhat\in\C^{E\times E}$。$n_v=12$ 时 $E=24$，不是把 fine spin 直接删成 $3\times3$ link。

定义 checkerboard
\[
p(x)=(x+y+z+t)\bmod2,\quad p=0\;(\mathrm{even}),\quad p=1\;(\mathrm{odd}).
\]
hopping 改变 parity，onsite 不改变。full coarse assets 保留 $\Lambda_c$ 全格；MATPC 输入才压缩一半。固定前三维时 parity compact 的时间坐标为
\[
t_{\rm full}=2t_{\rm half}+[(p-x-y-z)\bmod2].
\]
\src{\pathref{cpp/cuda/qcu/src/apply_multigrid_strict.cu:108--159} 定义 full/half site；:230--282 使用 parity decode；\pathref{pyqcu/tools/_strict_galerkin.py:378--405} 检查周期 checkerboard 的偶数 extent。}

\section{聚合基、限制延拓和 Galerkin}
令 $V_{s_fc_f,s_cj}(x)$ 为 aggregate 内的正交基：
\begin{align}
(P\phi)_{s_fc_f}(x)&=\sum_{s_c,j}V_{s_fc_f,s_cj}(x)\phi_{s_cj}(X(x)),\\
(R\psi)_{s_cj}(X)&=\sum_{x\in\mathcal A_X}\sum_{s_f,c_f}
 V_{s_fc_f,s_cj}(x)^*\psi_{s_fc_f}(x),\qquad R=P^\dagger.
\end{align}
局部 CGS 给出
\[
V_X^\dagger V_X=I_E,\quad P^\dagger P=I_{\rm coarse},\quad
PP^\dagger\ne I_{\rm fine}.
\]
任意有效 action 的下一层为
\[
D_{\ell+1}=R_\ell A_\ell^{\rm eff}P_\ell,\qquad
A_\ell^{\rm eff}=D_\ell\ {\rm(raw)},\quad
A_\ell^{\rm eff}=X_\ell^{-1}D_\ell\ {\rm(strict\ PC)}.
\]
one-hop support 只允许 $\mathcal N(X)=\{X,X\pm\hat\mu\}$；发现更远 displacement 就不能静默压缩成 $X/Y$。
\src{QUDA transfer：\pathref{refer/git-rep/quda/lib/transfer.cpp:152--163}；PyQCU batched support：\pathref{pyqcu/tools/_strict_galerkin.py:489--720,678--699}。}

\section{field 分配、入口和 setup 工作区}
\code{DiracClover::createCoarseOp} 调 \code{CoarseOp}，\code{DiracCloverPC::createCoarseOp} 传入 PC 类型；源码在
\pathref{refer/git-rep/quda/lib/dirac_clover.cpp:95--103,263--271}。coarse field 分配在 \code{DiracCoarse::createY}：
\[
\code{nColor}=N_s^cn_v,\quad\code{link\_type}=COARSE\_LINKS,\quad
\code{siteSubset}=FULL,\quad\code{nFace}=1.
\]
\code{createYhat} 再分配 $\Yhat$ 和无 ghost 的 $X^{-1}$，见
\pathref{refer/git-rep/quda/lib/dirac_coarse.cpp:121--220}。GPU native order 与 CPU QDP order 分开；Y/X order 不一致会在 :474 或 :644 直接报错。

\code{calculateY} 计算 aggregation、spin block、face 数和 \code{bidirectional\_links}，见
\pathref{refer/git-rep/quda/lib/coarse_op.cuh:974--1058}。工作区是
\[
AV=C^{-1}V\quad\text{(Clover-PC)},\qquad UV=U\,V\quad\text{或}\quad U\,AV.
\]
先生成 UV，再做 $(AV)^\dagger UV$ 的 contraction。Clover inverse 的 packed chiral kernel 为
\pathref{refer/git-rep/quda/include/kernels/coarse_op_kernel.cuh:693--760}；非预条件 path 令 $AV=V$。

\section{$X/Y$ 的数学构造}
forward link 为
\[
Y_\mu^{\rm f}(X)=-\kappa\sum_{x\in\mathcal A_X}
(AV_X(x))^\dagger U_\mu(x)V_{X+\hat\mu}(x+\hat\mu).
\]
backward storage block 为
\[
Y_\mu^{\rm b}(X-\hat\mu)=-\kappa\sum_{x\in\mathcal A_X}
V_X(x)^\dagger U_\mu^\dagger(x-\hat\mu)AV_{X-\hat\mu}(x-\hat\mu).
\]
实际代码可把 $-\kappa$ 延后到 dslash 的 \code{dim\_collapse}，所以比较 buffer 时必须记录系数是否已吸收。\code{computeUV} 以 \code{mma\_nn} 做 $UV$，\code{multiplyVUV} 以 \code{mma\_tn} 做共轭转置左乘，见
\pathref{refer/git-rep/quda/include/kernels/coarse_op_kernel.cuh:255--315,1157--1226}。多个 fine site 和 basis pair 并发写同一粗块，故使用 complex atomic add。

full Clover onsite 是
\[
X_C(X)=\sum_{x\in\mathcal A_X}V_X(x)^\dagger C(x)V_X(x),
\]
\code{compute\_coarse\_clover} 在
\pathref{refer/git-rep/quda/include/kernels/coarse_op_kernel.cuh:1768--1836}；Wilson-like 分支还执行 $X\leftarrow X+I_E$，见 :1878--1898。coarse-to-coarse 重新投影当前层 $X_\ell$，不是读取 fine $C$。Clover-like 表示 dense local block 和 inverse，不表示 fine 的 packed Clover tensor。

Clover-PC、Coarse-PC、twisted-PC 或 \code{need\_bidirectional} 保存两向；普通 operator 可调用 \code{reverse}。diagonal coarse-spin block 保持符号，off-diagonal 改变符号，见
\pathref{refer/git-rep/quda/include/kernels/coarse_op_kernel.cuh:1840--1875}。

\section{storage ABI、dslash 和 $\Yhat$ 次序}
QUDA coarse kernel 注释给出
\[
/out(x)=\sum_\mu Y_{-\mu}(x)in(x+\mu)+Y_\mu^\dagger(x-\mu)in(x-\mu).
\]
forward 读源点 forward slot；backward 用 $x-\mu$ 的 storage block 并取 dagger，源码见
\pathref{refer/git-rep/quda/include/kernels/dslash_coarse.cuh:126--249}。

令 action 目标点为 $q$。backward 项为
\[
Y_\mu^{\rm b}(q-\hat\mu)^\dagger z(q-\hat\mu).
\]
左预条件后的项必须满足
\[
X^{-1}(q)Y_\mu^{\rm b}(q-\hat\mu)^\dagger
=\left[Y_\mu^{\rm b}(q-\hat\mu)X^{-\dagger}(q)\right]^\dagger,
\]
故
\[
\boxed{\Yhat_\mu^{\rm b}(q-\hat\mu)
=Y_\mu^{\rm b}(q-\hat\mu)X^{-\dagger}(q)}.
\]
错误的 $X^{-1}(q-\hat\mu)Y_\mu^{\rm b}(q-\hat\mu)$ 同时错 site、错侧；非对易 dense block 会暴露该错误。

\code{calculateYhat} 的真实顺序是：batch invert $X$；双向交换 $Y$ ghost；backward 计算 $YX^{-\dagger}$；forward 计算 $X^{-1}Y$；注入 backward bulk；只交换 forward ghost。源码注释和实现为
\pathref{refer/git-rep/quda/lib/coarse_op_preconditioned.in.cu:148--282}。不变量为
\[
X^{-1}X=I,\quad\Yhat^{\rm f}=X^{-1}Y^{\rm f},\quad
(\Yhat^{\rm b})^\dagger=X^{-1}(Y^{\rm b})^\dagger.
\]

\section{三种 dslash 语义及 Schur}
定义 raw hopping
\[
H_Yz(X)=\sum_\mu[Y_\mu^{\rm f}(X)z(X+\hat\mu)
+Y_\mu^{\rm b}(X-\hat\mu)^\dagger z(X-\hat\mu)].
\]
若 $-\kappa$ 未吸收进 link，$D_cz=Xz-\kappa H_Yz$。\code{DiracCoarse::Dslash} 传 \code{clover=false}；\code{M} 以 \code{clover=true,dslash=true,xpay=true} 组合，见
\pathref{refer/git-rep/quda/lib/dirac_coarse.cpp:387--433}。onsite kernel 是 $out+=X\,in$ 或 $X^\dagger in$，见
\pathref{refer/git-rep/quda/include/kernels/dslash_coarse.cuh:257--289}。

\code{DiracCoarsePC::Dslash} 选择 $\Yhat$ 且 \code{clover=false}，见
\pathref{refer/git-rep/quda/lib/dirac_coarse.cpp:506--519}；$X$ 只供 prepare/reconstruct/raw path，不能在 PC dslash 再乘一次。对 target parity：
\[
u_q=\widehat H_{qp}x_p,\qquad
v_p=\widehat H_{pq}u_q,\qquad
M_p^cx_p=x_p-v_p.
\]
:530--561 的 even-even/odd-odd 分支拒绝 full field；symmetric branch 两次 hop 后以 \code{xpay(-1)} 形成 $I-\widehat H\widehat H$。

按 parity 排列
\[
D=\begin{pmatrix}A_p&D_{pq}\\D_{qp}&A_q\end{pmatrix},
\quad x_q=A_q^{-1}(b_q-D_{qp}x_p),
\]
\[
S_p=A_p-D_{pq}A_q^{-1}D_{qp},\quad
M_p=A_p^{-1}S_p,\quad
b_p^{\rm prep}=A_p^{-1}(b_p-D_{pq}A_q^{-1}b_q).
\]
full solution 的 coarse prepare/reconstruct 为
\[
b_p^{\rm prep}=X_p^{-1}b_p-(X_p^{-1}D_{pq})X_q^{-1}b_q,\qquad
x_q=X_q^{-1}b_q-(X_q^{-1}D_{qp})x_p.
\]
源码先 inverse、再 raw dslash、最后 axpy，见
\pathref{refer/git-rep/quda/lib/dirac_coarse.cpp:570--625}。

\begin{longtable}{@{}L{.19\textwidth}L{.25\textwidth}L{.27\textwidth}L{.19\textwidth}@{}}
\caption{full、Clover-PC、coarse-PC 对照}\label{tab:variants}\\\toprule
输入&中间量&最终对象&双向\\\midrule\endfirsthead\toprule{}输入&中间量&最终对象&双向\\\midrule\endhead
Clover full&$AV=V$&$Y=V^\dagger UV$，$X=V^\dagger CV+I$（按 convention）&可 reverse\\
Clover-PC&$AV=C^{-1}V$&$Y$ 含 inverse 作用，MATPC 显式归一化&必须\\
coarse Dirac&当前层 $X_\ell,Y_\ell$&$X_{\ell+1}=V^\dagger X_\ell V$&传递\\
coarse PC&输入是 $\Yhat_\ell$&下一层 coarsen $\Yhat$，不是 raw $Y$&必须\\\bottomrule
\end{longtable}
\src{\pathref{refer/git-rep/quda/lib/dirac_coarse.cpp:628--647} 明确 coarse-PC coarsens $\Yhat$ rather than $Y$。}

\section{一段贯穿全流程的超长公式型伪代码}
以下一个连续 \code{longtable} 每行只表达一个动作、等式或分支，行号不拆成独立算法。

\begin{longtable}{@{}L{.96\textwidth}@{}}
\caption{从 null vectors 到递归求解的连续算法}\label{tab:longalgo}\\\toprule
\endfirsthead\toprule
\multicolumn{1}{l}{\small\itshape（续表：同一算法，行号连续）}\\\midrule\endhead\bottomrule\endfoot
\textbf{输入：} $\Lambda_f,U_\mu,C,V,b,\kappa,m,\mu$，operator type，target parity $p$。\\
\textbf{检查：} $e=12$，$E=2n_v$，block size 整除 fine extent，coarse extent 为正整数。\\
\textbf{检查：} periodic checkerboard 的每个非平凡 coarse extent 为偶数。\\
\textbf{定义：} $X(x)=\lfloor x/b\rfloor$，$\mathcal A_X=\{x:X(x)=X\}$。\\
\textbf{正交：} 每个 aggregate 内 null vectors 做两遍 block CGS。\\
\textbf{不变量：} $\norm{V_X^\dagger V_X-I_E}\le\varepsilon_{\rm orth}$，失败停止。\\
\textbf{定义：} $(P\phi)(x)=V_X\phi(X)$，$(R\psi)(X)=V_X^\dagger\psi|_{\mathcal A_X}$。\\
\textbf{Galerkin：} $D_c=R A_{\rm eff}P$；raw 取 $A_{\rm eff}=D_f$，strict-PC 取 $A_{\rm eff}=X_f^{-1}D_f$。\\
\textbf{支持：} 只保留 $\Delta=0,\pm\hat\mu$，更远项必须报错。\\
\textbf{分配：} full $Y$ 八方向 slot、full $X$、atomic $Y/X$ 和工作区 $UV/AV$。\\
\textbf{分配：} PC 再分配 full $\Yhat$ 与 scalar $X^{-1}$。\\
\textbf{精度：} half/quarter 使用 single accumulator，保存前估计 maximum 和 scale。\\
\textbf{交换：} 计算跨 rank 的 $UV$ 前交换 $V$ 的 one-face ghost。\\
\textbf{PC 分支：} $AV(x)\leftarrow C^{-1}(x)V(x)$；非 PC 令 $AV\leftarrow V$。\\
\textbf{方向：} 对 $d=0,1,2,3$ 设置 FORWARDS。\\
\textbf{forward UV：} $UV_d(x)\leftarrow U_d(x)V(x+\hat d)$。\\
\textbf{forward contraction：} $Y_d^{\rm f}(X)\mathrel{+}\leftarrow(AV_X)^\dagger UV_d$。\\
\textbf{forward coefficient：} 若未吸收系数，$Y_d^{\rm f}\leftarrow-\kappa Y_d^{\rm f}$。\\
\textbf{forward storage：} 写入 source aggregate $X$ 的 forward slot。\\
\textbf{backward UV：} $UV_d^{\rm b}(x)\leftarrow U_d^\dagger(x-\hat d)AV(x-\hat d)$。\\
\textbf{backward contraction：} $Y_d^{\rm b}(X-\hat d)\mathrel{+}\leftarrow V_X^\dagger UV_d^{\rm b}$。\\
\textbf{backward storage：} 写入 $X-\hat d$，action 端读取其 dagger。\\
\textbf{双向：} Clover-PC、Coarse-PC 或 \code{need\_bidirectional} 时保存两向。\\
\textbf{reverse：} 单向普通 operator 用 spin-projector 符号规则生成另一向。\\
\textbf{atomic：} 对 fine site、basis pair、parity 并发 complex atomic add。\\
\textbf{edge：} 同 aggregate 贡献进入 $X$，跨 aggregate 贡献进入 $Y$.\\
\textbf{Clover onsite：} $X_C\leftarrow V^\dagger C V$。\\
\textbf{identity onsite：} $X\leftarrow X+I_E$。\\
\textbf{coarse onsite：} $X_{\ell+1}\leftarrow V^\dagger X_\ell V$。\\
\textbf{scale：} 检查每个方向 maximum；超出旧 scale 时重放 rescale。\\
\textbf{convert：} atomic fixed-point buffer 解码为应用 field。\\
\textbf{raw action：} $D_cz=Xz-\kappa\sum_d(Y_d^{\rm f}z_{+d}+Y_d^{\rm b\dagger}z_{-d})$。\\
\textbf{raw forward：} 读取 $Y_d^{\rm f}(X)$ 和 $z(X+\hat d)$。\\
\textbf{raw backward：} 读取 $Y_d^{\rm b}(X-\hat d)^\dagger$ 和 $z(X-\hat d)$。\\
\textbf{raw onsite：} \code{applyClover} 计算 $out\mathrel{+}=X\,in$。\\
\textbf{raw M：} \code{clover=true,dslash=true,xpay=true}。\\
\textbf{求逆：} 对每个 coarse site batch invert $X$。\\
\textbf{Yhat forward：} $\Yhat_d^{\rm f}(X)\leftarrow X^{-1}(X)Y_d^{\rm f}(X)$。\\
\textbf{Yhat backward：} $\Yhat_d^{\rm b}(X-\hat d)\leftarrow Y_d^{\rm b}(X-\hat d)X^{-\dagger}(X)$。\\
\textbf{不变量：} $(\Yhat_d^{\rm b})^\dagger=X^{-1}(Y_d^{\rm b})^\dagger$。\\
\textbf{ghost：} backward inject 到上一节点，再只 exchange forward。\\
\textbf{PC hop：} $\widehat Hz=\sum_d(\Yhat_d^{\rm f}z_{+d}+\Yhat_d^{\rm b\dagger}z_{-d})$。\\
\textbf{MATPC 输入：} 必须是 compact target parity；full field 直接拒绝。\\
\textbf{MATPC 第一步：} $u_q\leftarrow\widehat H_{qp}x_p$。\\
\textbf{MATPC 第二步：} $v_p\leftarrow\widehat H_{pq}u_q$。\\
\textbf{MATPC 输出：} $M_p^cx_p\leftarrow x_p-v_p$。\\
\textbf{prepare：} $b_p^{\rm prep}\leftarrow X_p^{-1}b_p-(X_p^{-1}D_{pq})X_q^{-1}b_q$。\\
\textbf{reconstruct：} $x_q\leftarrow X_q^{-1}b_q-(X_q^{-1}D_{qp})x_p$。\\
\textbf{递归：} coarse-PC 下一层输入必须是 $\Yhat$，不是 raw $Y$。\\
\textbf{资产：} 缓存 $V,X,X^{-1},Y^{\rm f/b},\Yhat^{\rm f/b}$，parity 只是 view。\\
\textbf{QCU ABI：} strict transition 固定输出 $V$/raw $Y$/ $\Yhat$ / $(X,X^{-1})$。\\
\textbf{QCU full hop：} forward link 与 $z_{+d}$、backward link dagger 与 $z_{-d}$。\\
\textbf{QCU parity hop：} 先 decode full coordinate，再取 compact neighbor。\\
\textbf{Galerkin test：} $\epsilon_G=\norm{D_cz-RD_fPz}/\max(1,\norm{RD_fPz})$。\\
\textbf{transfer test：} $\epsilon_{PR}=|\ip{Pz}{f}-\ip{z}{Rf}|/\max(1,\cdot)$。\\
\textbf{inverse test：} $\epsilon_X=\norm{X^{-1}X-I}/\max(1,\norm I)$，非正规 $X$ 同测 $XX^{-1}$。\\
\textbf{direction test：} unit basis 分别比较 forward/backward storage。\\
\textbf{极限：} $X=I\Rightarrow\Yhat=Y$；$Y=0\Rightarrow$ hop 为零；$E=1$ 检查 site/parity。\\
\textbf{非对易测试：} $X=\begin{psmallmatrix}2&1\\0&3\end{psmallmatrix}$、$Y=\begin{psmallmatrix}1&0\\1&2\end{psmallmatrix}$ 时 $X^{-1}Y\ne YX^{-1}$。\\
\textbf{MPI：} 单 rank 过 bulk，再在 boundary 比较 ghost 前后的 action。\\
\textbf{精度：} double reference、single compute、half storage 分别报告 inverse/Galerkin/solve residual。\\
\textbf{停止：} support、transfer、inverse、storage、ghost、精度不变量全通过才进入下一层。\\
\textbf{递归求解：} 以当前 coarse operator 为下一层 fine operator，重复 basis、Galerkin、$X/Y$、可选 $\Yhat$。\\
\textbf{输出：} field、geometry、$E$、precision、scale、parity metadata 和验证报告。\\
\end{longtable}
\src{QUDA 源：\pathref{refer/git-rep/quda/lib/coarse_op.cuh:974--1463}、\pathref{refer/git-rep/quda/lib/coarse_op_preconditioned.in.cu:148--282}、\pathref{refer/git-rep/quda/lib/dirac_coarse.cpp:506--647}；PyQCU：\pathref{pyqcu/tools/_strict_galerkin.py:334--720}。}

\section{递归、PyQCU 对照和验证}
每层资产为
\[
\mathcal G_\ell=\{V_\ell,X_{\ell+1},X_{\ell+1}^{-1},
Y_{\ell+1}^{\rm f/b},\Yhat_{\ell+1}^{\rm f/b}\}.
\]
V-cycle 递推为
\[
r_\ell=b_\ell-A_\ell x_\ell,\quad r_{\ell+1}=R_\ell r_\ell,\quad
e_{\ell+1}\simeq A_{\ell+1}^{-1}r_{\ell+1},\quad
x_\ell\leftarrow x_\ell+P_\ell e_{\ell+1}.
\]
PyQCU strict 的 \code{\_strict\_v\_cycle} 在
\pathref{pyqcu/solver/_quda_multigrid.py:3005--3048} 先在 target parity 平滑，再恢复 full coarse rhs；\code{diagnostics} 的 Galerkin 对照为 :3396--3402。QCU full kernel 在
\pathref{cpp/cuda/qcu/src/apply_multigrid_strict.cu:161--282} 与 QUDA storage 同构。

验证量为
\begin{align}
\epsilon_{PR}&=\frac{|\ip{Pz}{f}-\ip{z}{Rf}|}
{\max(1,|\ip{Pz}{f}|,|\ip{z}{Rf}|)},&
\epsilon_X&=\frac{\norm{X^{-1}X-I}}{\max(1,\norm I)},\\
\epsilon_G&=\frac{\norm{D_cz-RD_fPz}}{\max(1,\norm{RD_fPz})},&
\epsilon_b&=\frac{\norm{(\Yhat^{\rm b})^\dagger-X^{-1}(Y^{\rm b})^\dagger}}
{\max(1,\norm{X^{-1}(Y^{\rm b})^\dagger})}.
\end{align}

\begin{longtable}{@{}L{.24\textwidth}L{.31\textwidth}L{.31\textwidth}@{}}
\caption{极限测试与错误定位}\label{tab:checks}\\\toprule
情形&预期&定位\\\midrule\endfirsthead\toprule{}情形&预期&定位\\\midrule\endhead
$X=I$&$X^{-1}=I,\Yhat=Y$&inverse、方向或 scale\\
$Y=0$&只剩 onsite，PC hop 为零&clover flag 或 xpay\\
$E=1$&左右乘退化为标量&site/parity 错位\\
单 block&off-block link 消失&aggregate edge\\
周期边界&零点 backward 读取末点 storage&ghost inject/exchange\\
非对易 $X,Y$&左右乘结果不同&$\Yhat$ 次序\\
full MATPC input&拒绝或先 compact view&ABI 混用\\
half storage&解码矩阵接近 single reference&scale\\\bottomrule
\end{longtable}

GPU 逐元素回归必须固定 gauge、null vectors、geometry、precision 和 parity，同时报告 raw $X/Y$、$X^{-1}$、$\Yhat$、dslash、MATPC 和 Galerkin residual；solver residual 不能替代这些局部证据。本文生成过程未运行完整 GPU 逐元素回归，故该项为“未验证”。

\section{源码索引和最终不变量}
\begin{longtable}{@{}L{.40\textwidth}L{.48\textwidth}@{}}
\caption{直接引用源}\label{tab:sources}\\\toprule
源文件&内容\\\midrule\endfirsthead\toprule{}源文件&内容\\\midrule\endhead
\pathref{refer/git-rep/quda/lib/dirac_coarse.cpp:121--220}&Y/X/Yhat/Xinv 分配、geometry、order、ghost\\
\pathref{refer/git-rep/quda/lib/dirac_coarse.cpp:420--478}&raw M、prepare/reconstruct、coarse setup\\
\pathref{refer/git-rep/quda/lib/dirac_coarse.cpp:506--647}&PC dslash、MATPC、prepare/reconstruct、递归 $\Yhat$\\
\pathref{refer/git-rep/quda/lib/coarse_op.cuh:974--1463}&UV、AV、双向、atomic、reverse、onsite、scale\\
\pathref{refer/git-rep/quda/include/kernels/coarse_op_kernel.cuh:255--315,619--760,1157--1226}&UV、$C^{-1}V$、$(AV)^\dagger UV$\\
\pathref{refer/git-rep/quda/include/kernels/coarse_op_kernel.cuh:1608--1647,1768--1898}&edge、parity、Clover、identity\\
\pathref{refer/git-rep/quda/include/kernels/dslash_coarse.cuh:126--289}&forward/backward gather、dagger、onsite\\
\pathref{refer/git-rep/quda/lib/coarse_op_preconditioned.in.cu:148--282}&batch inverse、$YX^{-\dagger}$、$X^{-1}Y$、ghost\\
\pathref{pyqcu/tools/_strict_galerkin.py:334--720}&strict assets、Galerkin、support、MATPC\\
\pathref{pyqcu/solver/_quda_multigrid.py:3193--3402}&asset export、V-cycle、diagnostics\\
\pathref{cpp/cuda/qcu/src/apply_multigrid_strict.cu:161--282}&QCU full/parity hop\\
\pathref{docs/report_multigrid_quda_pyqcu_20260909.tex:399--691}&原报告 Schur、P/R、Galerkin、$X/Y/\Yhat$ 记号\\\bottomrule
\end{longtable}

\begin{enumerate}[leftmargin=2.5em]
\item 每层 $E=2n_v$；$X,Y,\Yhat$ 是 dense coarse blocks，不是 $SU(3)$。
\item raw backward 必须取 storage link dagger。
\item $\Yhat^{\rm f}=X^{-1}Y^{\rm f}$，$\Yhat^{\rm b}=Y^{\rm b}X^{-\dagger}$。
\item PC dslash 只做 $\Yhat$ hopping；MATPC 是 $I-\widehat H\widehat H$。
\item Clover-PC/coarse-PC 双向构造；递归 coarse-PC coarsen $\Yhat$。
\item full assets 与 compact parity view 分离。
\item 验证必须覆盖 transfer、inverse、Galerkin、storage、ghost 和精度。
\end{enumerate}
\src{源码与公式解析已完成；完整 GPU 逐元素回归依赖本机 CUDA/QUDA 构建环境，本次未宣称已完成。}
\end{document}
